\chapter{系统测试与分析}

本章旨在通过严谨的实验设计与多维度的性能评估，对本文提出的面向AI多任务场景的资源调度系统进行全面系统性的验证。实验依托腾讯TEG真实业务场景构建测试基准，重点围绕API接入层的适配鲁棒性、任务调度层的决策自动化程度以及任务执行层的资源互斥治理效能展开。通过对系统各功能模块及非功能性指标的深度压测，本章不仅验证了分层解耦架构在解决异构数据纠偏与算力冲突难题中的有效性，更进一步论证了系统在提升大规模AI任务并发周转率及优化算力资源周转效能方面的显著优势，为后续系统的工业级应用与迭代优化提供了详实的数据支撑。

\section{测试概述}

\subsection{开发与测试环境}

本系统的开发环境构建于高性能硬件基础之上，采用 Intel Core i9-13900K 处理器配合 64GB DDR5 高速内存，以满足调度系统在多容器并发运行及复杂逻辑编译时的算力需求。特别地，配置 NVIDIA GeForce RTX 4090 (24GB 显存) 显卡，旨在为 AI 大模型任务的参数画像识别与本地逻辑调试提供必要的算力支撑。软件层面，系统以 Debian 为核心开发系统，利用 Python 3.10 实现 API 接入层的数据纠偏与调度算法的核心逻辑。通过 Docker 容器化技术，实现了开发环境与运行环境的高度隔离，确保了异构数据映射机制在不同计算节点间的可移植性与一致性。本系统的开发环境如表6.1所示。

\begin{table}[htbp]
\centering
\caption{系统开发环境表}
\begin{tabularx}{\textwidth}{llXl}
\toprule
\textbf{环境类别} & \textbf{配置项目} & \textbf{详细规格} & \textbf{主要用途} \\
\midrule
硬件资源 & CPU & Intel Core i9-13900K @ 3.0GHz & 逻辑编译与计算 \\
硬件资源 & 内存 & 64GB DDR5 5600MHz & 多容器并行运行 \\
硬件资源 & GPU & NVIDIA GeForce RTX 4090 & 本地 AI 模型调试 \\
操作系统 & 开发系统 & Debian & 核心开发环境 \\
编程语言 & Python & Version 3.10 & 核心后端逻辑实现 \\
开发工具 & IDE & VS Code & 代码编写与断点调试 \\
容器化 & Docker & Version 24.0.5 & 开发环境隔离与镜像构建 \\
\bottomrule
\end{tabularx}
\end{table}

为深度验证系统在 AI 多任务场景下的调度性能与资源隔离能力，本研究搭建了一套基于 Kubernetes (Tencent TKE) 的分布式集群测试环境，由 1 台主节点与 2 台工作节点组成，模拟真实工业生产中的分布式计算场景。测试流程利用 Postman 对 API 接入层的鉴权与纠错功能进行全量覆盖验证，并借助 Hey 高性能压测工具模拟高并发任务请求，以此评估系统在极限负载下的吞吐量（RPS）与任务调度延迟。此外，环境涵盖了 CentOS、Debian 及 Windows Server 等异构操作系统，重点验证系统对不同计算底座的兼容性适配。所有测试数据通过 Grafana 进行可视化监控与分析，确保资源互斥智能算法在处理高价值算力冲突时的稳定性和鲁棒性，最终通过 腾讯智研 CI/CD 平台实现研发效能的全流程自动化验证。系统测试与验证环境如表6.2所示。

\begin{table}[htbp]
\centering
\caption{系统测试与验证环境表}
\begin{tabularx}{\textwidth}{llXl}
\toprule
\textbf{测试维度} & \textbf{测试项} & \textbf{环境配置/测试工具} & \textbf{测试目标} \\
\midrule
集群模拟 & 节点数量 & 3台分布式计算节点 (Master*1, Worker*2) & 验证分布式调度能力 \\
编排系统 & Kubernetes & Version 1.28.0 (Tencent TKE 托管版) & 资源动态扩缩容测试 \\
接口测试 & API 调试 & Postman & 后端接口功能验证 \\
性能压测 & 并发测试 & Hey (Golang 编写) & 评估系统吞吐量 (RPS) \\
可视化 & 面板展示 & Grafana & 生成稳定性压力测试报表 \\
兼容性 & 操作系统 & CentOS 7.9, Debian 11, Windows Server & 验证异构环境适配性 \\
自动化测试 & CI/CD 流程 & 腾讯智研 & 验证代码迭代自动化测试 \\
\bottomrule
\end{tabularx}
\end{table}

\subsection{测试方法}

本研究构建了多维度、全链路的测试体系，旨在验证系统在解决 AI 特性适配性、调度决策自动化以及计算资源互斥处理三大核心问题上的有效性。测试工作主要分为功能性测试与非功能性测试。

功能性测试通过开发自测（Unit Testing）与QA团队全案测试相结合，重点校验以下三个层面的逻辑正确性：AI适配与数据纠偏测试利用Postman构造包含全角字符、冗余符号、类型偏差的异构参数报文，验证API接入层在大模型任务参数画像约束下的自动纠正能力，确保下行至执行层的数据纯净度；决策引擎与可用性检测验证模拟不同系统负载水位与运营策略变化，验证调度层逻辑引擎是否能自主判定任务准入状态，消除对人工干预的依赖；资源互斥与并发调度验证通过构造多组竞争同一GPU编号的互斥任务流，验证系统是否能精准识别资源依赖拓扑，在保障临界资源不冲突的前提下，实现由串行向安全并发的转变。

非功能性测试侧重于系统在生产环境（腾讯TEG业务场景）下的工业级表现：性能压力测试采用基于Golang的Hey工具模拟高频并发请求，评估系统在多任务注入时的响应延迟与吞吐量（RPS），确保调度算法在毫秒级内完成决策；稳定性与容错测试依托公司基础设施进行7x24小时长周期运行监测，模拟节点宕机、网络抖动及内存受限（OOM）等异常场景，验证任务状态机的自愈能力与资源自动释放机制。异构兼容性测试在TKE托管的异构计算集群中，验证系统对CentOS、Debian及Windows Server等操作系统的适配性，确保参数传递机制在分布式节点间的一致性。

系统采用敏捷迭代模式，通过TAPD平台实现缺陷的生命周期管理。开发人员完成初测后，由QA团队进行全方位的回归测试。程序上线后，通过自动化监控系统捕获实时性能指标，并设立操作员反馈渠道，针对实际运行中出现的特殊AI算子适配需求或性能瓶颈进行持续优化与策略更新。

\section{功能测试}

功能测试部分通过对API接入层、任务调度层及任务执行层的全方位测试，系统性验证了本文所设计的资源调度系统在应对AI多任务复杂需求时的有效性。测试过程紧扣本文研究重点，重点针对异构数据适配、调度决策自动化及算力资源互斥治理三大核心问题进行了闭环验证：在API接入层测试中，验证了基于参数画像的自动化纠偏逻辑对AI任务异构数据的鲁棒性支持；在任务调度层测试中，证实了环境感知驱动的逻辑引擎能够准确判定任务可执行性，实现了调度流程的去人工化；在任务执行层与资源管理测试中，验证了资源互斥智能分析算法对GPU等独占性资源的高效管控，解决了因资源抢占导致的任务串行瓶颈。实验数据表明，系统各项核心指标均达到预期设计要求，功能用例通过率为100\%，为实现AI多任务的高效并发调度提供了坚实的技术保障。

\subsection{API接入层测试}

本次 API 功能接入层测试围绕权限鉴定与数据校验两大模块展开，模拟多场景验证功能准确性，所有案例结果均符合预期。权限鉴定模块中，组内用户登录可顺利完成身份验证，组外用户登录被拒绝；组内用户登录后能成功发起任务，组外用户无论是否登录，发起任务均被阻断并提示无权限，有效划分权限边界；数据校验模块里，提交合法数据可正常流转至业务模块；提交语法错误数据（如分隔符、引号问题），接入层会自动更正后继续处理；提交类型错误数据（如年龄字段为字符串、日期格式不符），则会阻断任务并通知操作员，避免错误数据影响下游业务。
API接入层数据校验与自动纠正功能测试用例如表6.3所示。

\begin{table}[htbp]
\centering
\caption{API接入层数据校验与自动纠正功能测试表}
\begin{tabularx}{\textwidth}{lXll}
\toprule
\textbf{需求项} & \textbf{预期结果} & \textbf{实际结果} & \textbf{测试状态} \\
\midrule
AI参数画像匹配 & 成功匹配Prompt/超参数Schema，系统判定为合法通过 & 100\%符合预期 & 通过 \\
异构类型校验   & 基于Pydantic实施类型检查，确保异构数据类型合法通过 & 100\%符合预期 & 通过 \\
全角标点纠偏   & 系统自动将全角符号转化为半角标点，实现修正通过 & 100\%符合预期 & 通过 \\
冗余符号清理   & 自动剔除输入中的冗余括号、空格等，实现修正通过 & 100\%符合预期 & 通过 \\
数值类型推断   & 实现浮点数向整数的无损转化，确保数据修正通过 & 100\%符合预期 & 通过 \\
字符串格式化   & 成功去除字符串首尾空格及不可见字符，实现修正通过 & 100\%符合预期 & 通过 \\
日期格式统合   & 各种日期字符串统一格式化为“YYYY-MM-DD”并修正通过 & 100\%符合预期 & 通过 \\
校验结果反馈   & 系统完整汇总校验报告并成功向调度平台下发提示 & 100\%符合预期 & 通过 \\
\bottomrule
\end{tabularx}
\label{tab:api_access_test}
\end{table}

此外，本工作还对权限鉴定、任务管理及脚本审核等核心基础功能进行了全面验证。权限鉴定测试结果显示组内与组外用户的拦截逻辑均100\%符合预期；任务管理测试结果显示本系统实现了任务表单的准确入库与状态机逻辑的正常流转；脚本审核功能验证了系统对AI执行逻辑安全性的严密管控，确保未获审批的脚本禁止执行。测试项的全部通过，标志着API接入层在安全准入与基础治理维度已具备支撑大规模异构AI任务高效流转的工业级可靠性。
API接入层其余测试用例表如表6.4所示。

\begin{table}[htbp]
\centering
\caption{API接入层其余功能测试表}
\begin{tabularx}{\textwidth}{lXll}
\toprule
\textbf{需求项} & \textbf{预期结果} & \textbf{实际结果} & \textbf{测试状态} \\
\midrule
权限鉴定 & 组内用户通过，组外用户拦截并提示无权限 & 100\%符合预期 & 通过 \\
任务管理 & 任务表单正确写入数据库，状态流转正常 & 100\%符合预期 & 通过 \\
脚本审核 & 审核通过后脚本上线，未审核脚本禁止执行 & 100\%符合预期 & 通过 \\
\bottomrule
\end{tabularx}
\end{table}

\subsection{任务调度层测试}

任务调度层测试聚焦任务可用性、资源可用性、容错与恢复三大核心维度，通过模拟不同业务场景验证调度功能的准确性与稳定性，所有测试案例执行结果均符合预期。

任务可用性测试中，对于过期任务，调度层能精准识别其时间属性并判定超出执行范围，不予调度；对于普通有效任务，调度层则按预设规则完成解析、优先级排序与指令下发，确保正确调度，从而保障任务执行的时间有效性。针对决策依赖难题，本测试验证了“场景策略注入”功能，模拟节假日或运维封网时，系统根据动态加载的JSON策略文件，自动调整任务准入优先级，展现了极高的灵活性与策略一致性。测试进一步验证了负载回落后的系统表现。当环境指标降至安全区间后，被挂起的任务能自动被调度器唤醒并重新入队执行，全程无需人工重启指令，实现了调度流程的闭环管理。任务可用性检测功能的 8 项核心用例通过率均为 100\%。实验数据论证了逻辑引擎在处理复杂环境约束时的鲁棒性，证实了系统已成功由“人工经验驱动”转型为基于实时状态感知的“智能自动化调度”，大幅降低了运维人力投入。任务可用性检测功能测试表如表 6.5 所示。

\begin{table}[htbp]
\centering
\caption{任务可用性检测功能测试表}
\begin{tabularx}{\textwidth}{lXll}
\toprule
\textbf{需求项} & \textbf{预期结果} & \textbf{实际结果} & \textbf{测试状态} \\
\midrule
系统负载检测 & 系统能够实时采集各项负载指标，并准确输出当前负载状态值 & 100\%符合预期 & 通过 \\
负载阈值配置 & 管理员可成功设置自定义负载阈值，系统保存后配置立即生效 & 100\%符合预期 & 通过 \\
任务启动判断 & 系统通过对比当前负载与预设阈值，准确作出允许或禁止启动的决策 & 100\%符合预期 & 通过 \\
高负载任务限制 & 当系统处于过载状态时，能有效拦截高耗能任务并弹出相应提示 & 100\%符合预期 & 通过 \\
配置修改记录 & 系统能完整记录管理员的配置修改操作，并成功生成审计日志 & 100\%符合预期 & 通过 \\
异常场景触发 & 在节假日等预设异常场景下，系统能自动激活相应的任务限制规则 & 100\%符合预期 & 通过 \\
任务拦截反馈 & 任务因过载被拦截后，系统能清晰向用户告知“负载过高”的具体原因 & 100\%符合预期 & 通过 \\
负载回落恢复 & 当实时负载回落至阈值以下时，系统能自动恢复任务的正常启动功能 & 100\%符合预期 & 通过 \\
\bottomrule
\end{tabularx}
\label{tab:task_scheduling_test}
\end{table}

针对人工智能多任务场景下的算力冲突痛点，资源可用性测试紧密围绕任务需求与底层计算资源的匹配关系展开。系统首先依托资源提前注册与任务资源声明功能，要求任务在提交阶段即标注具体的临界资源需求，从而确保算力元数据能够准确入库并建立稳固的任务关联。在核心的调度决策阶段，调度层深度利用任务互斥判定机制，对在途任务的临界资源需求进行实时对比分析。若判定目标资源处于空闲状态，则通过执行权限判断模块立即下发启动指令；而一旦检测到资源占用冲突，系统则会自动拦截任务下发并向操作员反馈清晰的互斥冲突提示，同时将该任务平滑地加入高可靠的资源等待队列中。此外，调度层内置的资源释放监测模块会全时段实时跟踪算力资源的动态占用变化。一旦捕捉到资源由锁定转为空闲的瞬时状态，系统便会立即触发等待任务唤醒机制，驱动执行队列进行自动调度与分发，从而有效缩短任务间的间隙空耗。在容错与恢复测试维度，系统通过模拟计算节点关机或断电等极端硬件故障场景，充分验证了调度层能凭借心跳检测技术敏锐识别节点的离线状态。系统随后会迅速触发预设的故障转移机制，将尚未完成的算力任务重新分配至集群内的其他可用节点执行，从底层有效保障了AI业务的连续性与系统运行的稳健性。本项资源可用性检测的具体功能测试用例与验证结果详见表6.6。

\begin{table}[htbp]
\centering
\caption{资源可用性检测功能测试表}
\begin{tabularx}{\textwidth}{lXll}
\toprule
\textbf{需求项} & \textbf{预期结果} & \textbf{实际结果} & \textbf{测试状态} \\
\midrule
资源提前注册 & 任务能够声明所需资源并完成登记，确保资源信息成功入库 & 100\%符合预期 & 通过 \\
任务资源声明 & 任务提交时可标注具体资源需求，系统成功记录任务与资源的关联关系 & 100\%符合预期 & 通过 \\
任务互斥判定 & 系统对比任务所需的临界资源，能够准确识别并标记存在互斥关系的任务 & 100\%符合预期 & 通过 \\
执行权限判断 & 检查目标资源的占用状态，并根据检测结果准确作出允许或拒绝执行的决策 & 100\%符合预期 & 通过 \\
任务等待入队 & 在无可用资源时自动标记任务状态，并将其平滑加入资源等待队列中 & 100\%符合预期 & 通过 \\
资源释放监测 & 实时跟踪并捕捉资源占用的动态变化，确保能够及时识别资源的空闲状态 & 100\%符合预期 & 通过 \\
等待任务唤醒 & 当资源释放并处于空闲状态时，成功触发队列调度并唤醒等待中的任务 & 100\%符合预期 & 通过 \\
互斥冲突提示 & 当发生任务互斥导致无法执行时，系统能清晰向用户提示资源占用原因 & 100\%符合预期 & 通过 \\
\bottomrule
\end{tabularx}
\label{tab:resource_availability_test}
\end{table}

针对任务调度层除了可用性检测以外的其他关键职能，本测试汇总了负载均衡、数据同步以及容错恢复维度的测试验证结果，这些测试项旨在验证系统在分布式异构环境下保障执行效能与运行韧性的能力，深度契合了第三章中提出的各项非核心功能需求。

在负载均衡与性能优化方面，测试重点针对任务执行节点在物理基础与运行环境上的显著异构性展开。测试证实，调度层通过引入智能权重分配算法，能精准采集各节点的实时负载指标并动态分析AI任务的资源亲和性特征，将计算密集型推理任务优先指派至算力充裕的GPU节点，从而在多任务高并发场景下维持高吞吐与低延迟的稳健运行状态。在数据同步维度，本系统构建了高可靠、强一致的同步机制以规避子模块间因信息流转失真导致的任务逻辑偏离。测试验证了基于高性能中间件的消息同步底座是否能实现数据的透明流转，确保API接入层预先完成的参数纠偏与类型推断结果能实时、完整地透传至任务执行节点，从而为算力的精准调度提供坚实的数据基础。在容错与恢复能力验证中，考虑到AI任务具备执行周期长且算力消耗高的特征，系统必须具备对非致命性异常的动态识别与自助修复能力。实验重点模拟了网络通信韧性、磁盘空间不足、以及因异构环境导致的参数类型冲突等典型异常场景，验证了系统是否能依据指数退避机制执行自动重试，并利用智能纠偏规则对float与int类型冲突进行实时修复。测试结果显示，系统在遭遇计算节点宕机或网络抖动时，能迅速通过状态机重置触发故障转移机制，将任务重新分配至可用节点恢复执行，成功实现了由人工排障向引擎自愈的深度转型，各项指标均100\%符合预期。任务调度层其余功能测试表如表6.7所示。

\begin{table}[htbp]
\centering
\caption{任务调度层其余功能测试表}
\begin{tabularx}{\textwidth}{lXll}
\toprule
\textbf{需求项} & \textbf{预期结果} & \textbf{实际结果} & \textbf{测试状态} \\
\midrule
负载均衡 & 任务按权重比例分发至异构执行节点 & 100\%符合预期 & 通过 \\
数据同步 & 任务池状态与执行队列实时保持一致 & 100\%符合预期 & 通过 \\
容错恢复 & 任务重新入队并在可用节点恢复执行 & 100\%符合预期 & 通过 \\
\bottomrule
\end{tabularx}
\end{table}

\subsection{任务执行层测试}

任务执行层作为算力调度的最终落地单元，其测试核心在于验证分布式环境下参数传递的准确性、任务动态编辑的即时性以及异常捕获机制的鲁棒性。本测试旨在确保调度中枢下发的逻辑指令能精准转化为底层的算力输出，并验证系统在极端执行工况下的自愈与反馈能力。

为了验证系统是否有效解决AI特性适配不佳及异构数据处理能力不足的痛点，本研究对参数传递环节进行了全方位的覆盖测试。系统通过“参数封装处理”与“传输校验机制”，确保了异构参数在分布式节点间流转时的完整性系统内置了“任务丢失检测”与“断点续传支持”，测试结果表明，系统在发现传输中断后能准确触发重试逻辑，并支持大参数的分片传输，有效避免了因网络异常导致的算力资源空耗。系统重点验证了“远程附件拉取”与“参数路径替换”功能。实验通过构造包含多组对象存储（COS）链接的报文，证实执行层能自动化完成远程素材的并发下载，并将参数中的远程URL动态重定向为本地物理路径。这充分论证了本文设计的参数传递协议与环境自愈逻辑具备极高的鲁棒性，为后续实现复杂的自动化决策与资源精细化治理提供了坚实的数据闭环保障。参数传递功能测试表如表6.8所示。

\begin{table}[htbp]
\centering
\caption{参数传递功能测试表}
\begin{tabularx}{\textwidth}{lXll}
\toprule
\textbf{需求项} & \textbf{预期结果} & \textbf{实际结果} & \textbf{测试状态} \\
\midrule
参数封装处理 & 能够打包任意类型参数并成功生成传输数据包  & 100\%符合预期 & 通过 \\
透明传输处理 & 成功屏蔽网络传输细节并实现无感调用  & 100\%符合预期 & 通过 \\
传输校验机制 & 准确验证参数完整性并确保数据准确  & 100\%符合预期 & 通过 \\
任务丢失检测 & 实时监控传输状态并在发现丢失时触发重试  & 100\%符合预期 & 通过 \\
远程附件拉取 & 成功下载远程附件至本地并生成本地文件路径  & 100\%符合预期 & 通过 \\
参数路径替换 & 将远程路径改为本地路径以适配本地执行环境  & 100\%符合预期 & 通过 \\
断点续传支持 & 实现大参数分片传输并支持中断后继续  & 100\%符合预期 & 通过 \\
传输结果反馈 & 准确返回参数处理状态并告知成功或失败原因  & 100\%符合预期 & 通过 \\
\bottomrule
\end{tabularx}
\end{table}

本测试证实了执行层在AI多任务管理中的灵活性与可靠性。系统能精准执行撤销指令并终止无效进程，通过即时释放GPU等独占资源规避算力空耗。此外，系统支持异构参数编辑，并能驱动新实例重新联动决策引擎与互斥分析逻辑。100\%的操作反馈准确率保障了任务全生命周期的观测透明度，相关数据见表6.9。

\begin{table}[htbp]
\centering
\caption{任务编辑功能测试表}
\begin{tabularx}{\textwidth}{lXll}
\toprule
\textbf{需求项} & \textbf{预期结果} & \textbf{实际结果} & \textbf{测试状态} \\
\midrule
任务撤销触发 & 能够正确接收撤销指令并启动任务终止流程 & 100\%符合预期 & 通过  \\
任务参数修改 & 成功完成任务参数及类型的编辑并保存修改内容 & 100\%符合预期 & 通过  \\
任务重启执行 & 在符合条件时成功启动新任务并进入运行态 & 100\%符合预期 & 通过  \\
操作结果反馈 & 准确返回撤销或重启状态并告知成功或失败原因 & 100\%符合预期 & 通过  \\
\bottomrule
\end{tabularx}
\end{table}

本工作针对人工智能大模型调度场景下任务高算力消耗且周期长的特征，对系统全链路的风险识别与自愈能力进行了严谨验证。测试结果证实，系统能够实时细致地监测CPU、内存及文件描述符等关键指标，并基于历史数据准确触发预警通知，从而在感知识别维度保障了运行环境的稳定性。在拦截与保护层面，异常任务暂停与数据暂存功能的通过，确保了系统能在终止风险任务的同时保存中间计算成果，从底层化解了因非预期中断导致的算力空耗难题。此外，依托Python异常处理流程实现的自动化日志捕获、信息即时上报及原因深度分析，显著提升了运维反馈的精确度，配合100\%符合预期的系统恢复尝试逻辑，标志着系统已成功实现从“人工排障”向“引擎自愈”的智能化转型。全体测试项的顺利通过，充分论证了执行层在复杂异构工况下具备极高的鲁棒性，为支撑大规模AI任务的高效并发周转构筑了稳固的安全防线。异常捕获功能测试表如表6.10所示。

\begin{table}[htbp]
\centering
\caption{异常捕获功能测试表}
\begin{tabularx}{\textwidth}{lXll}
\toprule
\textbf{需求项} & \textbf{预期结果} & \textbf{实际结果} & \textbf{测试状态} \\
\midrule
系统指标监测 & 能够实时、细致地监测CPU、内存及文件描述符等关键指标并输出数据 & 100\%符合预期 & 通过  \\
异常预测预警 & 基于历史数据准确识别指标异常，并成功向相关人员触发预警通知 & 100\%符合预期 & 通过  \\
异常任务暂停 & 发现系统风险时能够即刻终止目标任务，避免不必要的资源浪费 & 100\%符合预期 & 通过  \\
任务数据暂存 & 在异常任务暂停或终止前，成功保存中间数据以防止工作量丢失 & 100\%符合预期 & 通过  \\
异常捕获 & 依托Python异常处理流程，成功采集任务运行时产生的完整异常日志 & 100\%符合预期 & 通过  \\
异常信息上报 & 将详尽的异常详情及日志信息即时推送至相关人员以通知故障处理 & 100\%符合预期 & 通过  \\
异常原因分析 & 自动整理并提供Traceback日志及数据报表，为故障排查提供分析依据 & 100\%符合预期 & 通过  \\
系统恢复尝试 & 能够通过重启服务、释放被占用的临界资源等手段恢复系统正常运行 & 100\%符合预期 & 通过  \\
\bottomrule
\end{tabularx}
\end{table}

\section{非功能测试}

\subsection{性能测试}

性能测试不仅是验证系统基础吞吐能力的手段，更是评估资源互斥智能分析与调度决策逻辑在极端环境下效能的关键环节。其核心意义在于通过高并发压力模拟，验证系统是否真正具备实时处理能力，从而确保系统能彻底消除传统串行调度模式引发的瓶颈，支撑AI算力资源的高效并发周转。性能测试的重点在于验证“资源互斥智能分析”与“自动化决策引擎”的响应延迟是否控制在极低阈值内。通过在高并发环境下实时采集调度指令的下发耗时，本测试能够论证系统在保障临界资源不冲突的前提下，实现由“保守串行”向“安全并发”转型的实际效能，进而为计算集群利用率的提升提供数据支撑。

为了精准匹配AI高频研发场景下对“无感官卡顿”的严苛需求，本测试选用了基于 Golang编写的高性能压测工具Hey。由于Golang在并发处理与运行性能上相较于Python具备显著优势\cite{wang2023}，能够更客观、高效地对系统后端实施极限压测，从而准确捕捉大规模附件流转及批量指令控制时的响应延迟。测试结果表明，系统在多任务注入场景下依然保持了毫秒级的调度响应速度，证实了其在处理AI异构任务时具备极高的鲁棒性与高吞吐处理能力，完全符合预期的性能目标。具体的性能数据如表6.11所示。

\begin{table}[htbp]
\centering
\caption{性能测试结果表}
\begin{tabular}{lll}
\toprule
测试内容 & 预期结果 & 实际结果 \\
\midrule
10 并发下提交任务平均响应时间 & 响应时间 < 50 ms & 平均 14 ms \\
10 并发下删除任务平均响应时间 & 响应时间 < 50 ms & 平均 11 ms \\
10 并发下编辑任务平均响应时间 & 响应时间 < 50 ms & 平均 13 ms \\
100 并发下提交任务平均响应时间 & 响应时间 < 100 ms & 平均 32 ms \\
100 并发下删除任务平均响应时间 & 响应时间 < 100 ms & 平均 28 ms \\
100 并发下编辑任务平均响应时间 & 响应时间 < 100 ms & 平均 30 ms \\
1000 并发下提交任务平均响应时间 & 响应时间 < 500 ms & 平均 168 ms \\
1000 并发下删除任务平均响应时间 & 响应时间 < 500 ms & 平均 145 ms \\
1000 并发下编辑任务平均响应时间 & 响应时间 < 500 ms & 平均 152 ms \\
\bottomrule
\end{tabular}
\label{tab:perf_concurrent}
\end{table}

\subsection{稳定性测试}

本工作通过对关机重启、断电、网络中断、多任务并发及磁盘异常等复杂工况的模拟，全面验证了系统在工业级生产环境下的长期运行可靠性与自动化修复能力。针对AI任务执行周期长且算力成本高的特征，测试证实系统在遭遇硬件故障或网络抖动时，能够依托环境感知的自愈机制实现程序的自动重连与未完成任务的断点续传，确保了数据完整性并有效规避了算力空耗。同时，多任务并发执行的成功校验进一步论证了资源互斥智能分析算法在处理高价值算力冲突时的精确度，保障了在共享资源环境下各任务结果的正确性。此外，系统在磁盘空间不足等极限场景下的安全停止与预警反馈，配合100\%符合预期的测试结果，标志着本研究已成功构建起覆盖全生命周期的稳健运行防线，从技术底层实现了调度决策的去人工化与业务连续性保障 。稳定性测试用例如表6.12所示。

\begin{table}[htbp]
\centering
\caption{稳定性测试用例表}
\begin{tabular}{lll}
\toprule
测试内容 & 预期结果 & 实际结果 \\
\midrule
关机重启 & 程序自动恢复，数据无丢失 &  100\%符合预期 \\
断电 & 继续未完成任务 &  100\%符合预期 \\
网络中断 & 程序重连网络，任务继续 &  100\%符合预期 \\
多任务并发执行 & 任务无冲突，结果正确 &  100\%符合预期 \\
磁盘空间不足时 & 提示空间不足，程序停止 &  100\%符合预期 \\
\bottomrule
\end{tabular}
\label{tab:my_label}
\end{table}

\subsection{兼容性测试}

兼容性测试旨在验证系统对多元化AI任务范式与多态运行环境的深度适配能力。测试重点评估了系统在异构计算集群中的鲁棒性，从根源上解决了子问题一中AI特性适配欠佳的难题。本工作针对AI任务强数据依赖的特征，测试了多类操作系统环境。验证证实，参数传递机制与远程附件拉取逻辑在跨集群分配策略下表现一致，确保了异构计算底座上的逻辑一致性与数据鲁棒性。系统在不同指令集架构与操作系统组合下均能稳定运行且无报错。测试结果表明，本文的设计能够应对环境偏差导致的适配风险，为系统的工业级推广提供了坚实的兼容性保障。验证结果如表6.13所示。

\begin {table}[htbp]
\centering
\caption {任务执行层兼容性测试用例表}
\begin {tabular}{lll}
\toprule
测试内容 & 预期结果 & 实际结果 \\
\midrule
Linux系统x86服务器 & 程序适配架构，运行稳定无报错 &  100\%符合预期 \\
Linux系统ARM服务器 & 程序适配架构，运行稳定无报错 &  100\%符合预期 \\
Windows系统x86服务器 & 程序适配架构，运行稳定无报错 &  100\%符合预期 \\
Windows系统ARM服务器 & 程序适配架构，运行稳定无报错 &  100\%符合预期 \\
\bottomrule
\end {tabular}
\label {tab:my_label}
\end {table}

\section{本章小结}

本章对面向AI多任务场景的资源调度系统进行了系统性的测试验证与结果分析。通过模拟真实的生产业务环境，从功能完备性、性能高效性、系统稳定性及环境兼容性四个维度开展了全方位评估。通过对异构数据纠偏、自动化准入决策及临界资源互斥识别等关键模块的验证，证实了本系统能有效解决AI任务参数复杂、调度依赖人工及计算资源冲突等行业痛点。

\makeatletter
\@openrightfalse
\makeatother